\documentclass[11pt]{article}
\usepackage[margin=1in]{geometry}
\usepackage{amsmath,amssymb}
\usepackage{booktabs}
\usepackage{hyperref}
\usepackage{enumitem}

\newcommand{\code}[1]{\texttt{#1}}

\title{v23: A Combined Architecture Specification\\
\large Synthesizing the v22 Baseline with the a100\_final Blur-Robustness Diagnosis}
\author{Internal architecture plan --- document only, no code}
\date{}

\begin{document}
\maketitle

\begin{abstract}
This document specifies \textbf{v23}, a single combined architecture that grafts the
demonstrated strengths of \code{train\_gvp\_egnn\_v22.py} (a strong, well-calibrated
zero-blur baseline) onto the diagnosed, evidence-based blur-robustness fixes developed in
the \code{a100\_final} rebuild. It is written to be buildable from this document alone: every
component names its exact source, its exact formula, and its exact integration point.
Two goals are pre-registered before any training happens, in the same spirit as every gate
used earlier in this project.
\end{abstract}

\tableofcontents

%==============================================================================
\section{Pre-registered goals}
%==============================================================================

\textbf{Goal 1 --- reconstruction.} Minimize median angular error and maximize head-tail
accuracy, beating the classical PCA/geometric baseline, in every regime where the physics
permits it. \emph{Scoped explicitly}: this session established, three independent ways
(an analytical statistical bound; the trained R7 network converging to 99.3\% numerical
agreement with plain PCA; a separate non-neural method using deliberately richer
hand-engineered shape features finding no rescue), that low energy (1--5~keV) under real
blur ($\sigma\gtrsim3$--10~nm) is a genuine, near-irreducible information floor for any
point-cloud-shape-based estimator. \textbf{Beating PCA at the low-energy/high-blur floor is
explicitly not a goal of v23.} The goal applies to: (a) near-zero blur across all three
energy tiers, and (b) mid/high energy across the full blur range $\sigma\in\{0,\dots,100\}$nm.

\textbf{Goal 2 --- calibration.} Genuine, well-calibrated uncertainty (empirical coverage
close to nominal at multiple confidence levels, low ECE) produced by the architecture, loss,
and training procedure themselves, \textbf{with no post-hoc temperature scaling or
recalibration map fit on held-out/eval data at any point.} Calibration is a first-class
design target here, not a bolt-on.

%==============================================================================
\section{Why v23 exists: the evidentiary chain}
%==============================================================================

\subsection{The winner lineage and its rebuild}
The project's original ``winner'' (\code{train\_diag\_r7\_trunc\_only.py}, the R7 lineage)
built direction estimation around a hard architectural constraint: the predicted direction
can \emph{only ever} be a re-weighted combination of a small pool of $\sim$12
mostly-parameter-free candidate vectors (the ``V-bottleneck''), specifically to guarantee
continuous SO(3) equivariance. \code{a100\_final} (full detail in its own
\code{README.md}) rebuilt this with five targeted fixes --- $\sigma$ fed into the
direction-weighting pathway, a blur-deconvolved radius-of-gyration normalizer, blur-corrected
higher-moment candidate vectors, a widened $K{=}10$ vMF mixture, and a blur-adaptive radius
graph. Measured result (the ``small'' config, matching the winner's capacity): genuine
improvement at low blur ($\sigma$=0--6~nm, e.g.\ high-E $\sigma$=0: 8.8$^\circ\to$5.7$^\circ$)
but \textbf{flat-to-worse at exactly the target regime} (e.g.\ high-E $\sigma$=50~nm:
32.8$^\circ\to$39.9$^\circ$, a 7.1$^\circ$ regression), and calibration \textbf{worse}, not
better (ECE $\sim$28--42\% vs.\ a $\sim$20\% prior baseline, \emph{worst specifically where
the model is most accurate/confident} --- high energy, low blur, ECE $\sim$40\%).

\subsection{The discovery of v22}
A completed, real run --- \code{python src/scripts/train\_gvp\_egnn\_v22.py
--directional-head --tag ROT\_FINAL} --- on the \emph{exact same} crystal-rotation corpus
(\code{data/siimpl\_rot/}) that all of the above used, achieves (from the real
\code{eval\_results.csv}, 49{,}200 rows / 1{,}000 configs, pooled across 1--105~keV, confirmed
zero-blur): \textbf{median angular error 6.8$^\circ$ (mean-readout) / 6.7$^\circ$
(mode-readout), head-tail 97.0\%/96.6\%, coverage gaps of only 2.6--3.3 percentage points}
(68\% nominal$\to$65.4\% actual, 90\% nominal$\to$86.7\% actual). This is dramatically better
on all three axes than anything the R7/\code{a100\_final} lineage produced, at zero blur.

\subsection{Why v22 is stronger --- ranked findings from a direct code comparison}
\begin{enumerate}[leftmargin=*]
\item \textbf{Head architecture (the largest factor).} v22's \code{DirectionalPosterior}
      (\code{src/models/directional\_head.py}) is an \emph{unconstrained} MLP
      (\code{VMFMixture.\_params}, lines 88--94: a 2-hidden-layer \code{\_MLP} with
      \code{hidden=256}) regressing $K{=}4$ vMF components' $(\mu,\kappa,\text{logit})$
      directly from the pooled conditioning vector $z$ --- no Gram-matrix step, no
      restriction to a small candidate-vector pool. Equivariance is approximated purely by
      the same 48-element discrete $O_h$ augmentation \code{a100\_final} also uses, not
      architecturally enforced. \textbf{Diamond's true symmetry is the discrete $O_h$
      crystal group, not continuous SO(3)} --- channeling is itself proof continuous
      rotational symmetry is broken. The R7/\code{a100\_final} lineage's architectural
      commitment to continuous equivariance (the V-bottleneck) may therefore have been
      enforcing an incorrect symmetry constraint the entire time, at a real expressiveness
      cost, for a guarantee the physics does not have.
\item \textbf{Two-stage training.} Stage 1 jointly trains the backbone $+$
      \code{DirectionalPosterior} $+$ two \emph{auxiliary} heads (a separate
      \code{DirectionHead} giving an auxiliary vMF loss, \code{EnergyHead} giving an
      auxiliary energy loss), with a sign-aware curriculum on the \emph{auxiliary} direction
      loss only (see \S\ref{sec:loss}). Stage 2 \textbf{freezes the entire backbone} and
      trains a fresh, newly-initialized \code{DirectionalPosterior} from scratch on
      Pool-A-only (isotropic, no channeling enrichment) data. The deployed/evaluated model is
      Stage 2's checkpoint.
\item \textbf{Physics-augmented scalar features.} v22 appends 5 pre-truncation scalar
      features to the conditioning vector (exact formula in \S\ref{sec:physfeat}).
      \code{a100\_final} has no equivalent.
\item $K{=}4$ mixture (v22) vs.\ $K{=}10$ (\code{a100\_final}) --- consistent with an older,
      independent hyperparameter sweep (frozen-backbone head sweep, prior session) that found
      widening past $K{=}4$ ``gave nothing,'' and consistent with \code{a100\_final}'s $K{=}10$
      correlating with \emph{worse}, not better, calibration.
\item The signed-loss curriculum (item 2, detailed in \S\ref{sec:loss}) plausibly explains
      v22's much higher head-tail accuracy (97\% vs.\ \code{a100\_final}'s $\sim$88--91\%,
      which uses a fully-decoupled separate sign-classifier).
\end{enumerate}

\subsection{v22 under blur: also broken, and diagnostically important}
\code{train\_gvp\_egnn\_dr\_sigma.py}'s own docstring records a zero-shot eval of the
deployed v22 checkpoint on blurred data collapsing from 6.7$^\circ$/96\% at $\sigma{=}0$ to
\textbf{77.8$^\circ$/57\% at just 1~nm}, saturating near-random by 3~nm --- a much
\emph{faster} collapse than R7's more gradual degradation, despite v22's far better starting
point. \textbf{Root cause, confirmed by reading \code{preprocess\_egnn} (\code{src/utils/data\_utils.py}):
v22's Stage 1 is called with the default \code{smear\_sigma=0.0} --- the backbone never sees
a single blurred training example.} Two fixes were already built and actually run to
completion:
\begin{itemize}[leftmargin=*]
\item \code{dr\_sigma}: warm-start from Stage-1 weights, then \textbf{unfreeze everything},
      fine-tune under blur (30 epochs). Real training log: $\sigma{=}0$ error \emph{drifted
      up} from 5.96$^\circ$ to $\sim$6.19$^\circ$.
\item \code{from\_scratch\_dr}: train from a fresh random init with \emph{continuous} blur
      baked into both stages from epoch 1 (300+21 epochs --- the most thorough possible blur
      exposure). Real training log: $\sigma{=}0$ error converged to $\sim$7.1--7.2$^\circ$.
\end{itemize}
The user who ran both confirms, directly: results under blur were \textbf{poor} for both.
Neither training log recorded per-sigma-tier held-out performance, so the exact failure
curve is not recoverable, but the qualitative outcome is confirmed by the person who ran it.

\subsection{The central argument for v23: three independent failures share one missing ingredient}
\label{sec:central-argument}
R7 (trained with dynamic blur throughout; this session established the mid/high-energy
degradation is \emph{bias}, not variance, via a noise-averaging test showing no benefit from
averaging predictions over repeated independent blur draws of the same track, replicated
twice with a hemisphere-alignment-bug control), \code{dr\_sigma} (backbone re-exposed to blur
via unfreezing), and \code{from\_scratch\_dr} (backbone exposed to the maximum possible blur
from initialization) all fail to achieve blur-robustness. \textbf{All three share exactly one
property: none of them ever feed the actual blur magnitude into the model as an explicit
input.} Blur is used only to generate noisy training data, never told to the network. Given
that ``expose the model to more/different blur'' has now failed three architecturally
distinct ways, the convergent explanation is that the model, regardless of how much blurred
data it has seen, has no way to know \emph{which} blur level the current input came from ---
so it can only learn one compromise behavior, averaged over the training blur distribution,
which is mathematically guaranteed to be systematically wrong (biased) for any specific
sigma. This is the single most important, best-evidenced new ingredient in v23.

\textbf{Honest caveat, stated plainly.} \code{a100\_final} \emph{did} include a form of
sigma-conditioning (feeding $\sigma$ into the V-bottleneck's coefficient-weighting pathway)
and its backbone \emph{was} already blur-exposed (inherited from R7's dynamic smearing), yet
the bundled result was still flat-to-worse at target cells. This was never isolated as a
single-variable test --- it was bundled with four other simultaneous changes, including the
$K{=}10$ widening independently suspected of hurting calibration, inside an architecture
(the V-bottleneck) that may have structurally limited how much use a flexible head could make
of the sigma signal even if it were informative. \textbf{v23's hypothesis is therefore not
identical to what \code{a100\_final} already tried}: it pairs sigma-conditioning with an
\emph{unconstrained} head (removing the V-bottleneck) and a backbone whose \emph{pretraining
stage itself} is exposed to blur (v22's Stage~1 currently is not, at all) --- a configuration
that has not actually been tested.

%==============================================================================
\section{v23 component specification}
%==============================================================================

\subsection{Backbone --- carried forward from v22, unchanged}
GVP-EGNN with \textbf{hidden\_dim=96, n\_layers=6, k\_neighbors=16 (replaced, see
\S\ref{sec:graph}), n\_heads=8, d\_proj=48, d\_latent=384, v\_dim=8, n\_axes=2}
(\code{train\_gvp\_egnn\_v22.py} lines 110--118). No evidence from this session suggests the
backbone family itself is the problem; the V-bottleneck removal is a \emph{head}-side change
(\S\ref{sec:head}). \textbf{Capacity is deliberately not re-opened as a question here}: a
prior-session finding (frozen-backbone head sweep, same lineage as the $K{=}4$ result cited
above) concluded this backbone is \emph{information-limited, not capacity-limited} at this
size, and that increasing width/depth wastes compute without recovering signal. v23 inherits
hidden=96 on that basis, not by default/oversight.

\subsection{Head --- v22's unconstrained design, with $\sigma$ added as a conditioning scalar}
\label{sec:head}
Keep \code{DirectionalPosterior} exactly as built: \code{VMFMixture} with
\textbf{$K{=}4$} components, a 2-hidden-layer MLP (\code{hidden=256}) mapping the
conditioning vector to $(\mu_k,\kappa_k,\text{logit}_k)_{k=1}^4$
(\code{directional\_head.py} lines 81--94), and \code{GMM1D} with $K{=}3$ components for
log-energy. \textbf{Do not reintroduce the V-bottleneck/candidate-vector-pool restriction.}

\textbf{Integration point for $\sigma$-conditioning, revised per external review to use a
sinusoidal embedding rather than raw scalar concatenation.} A fresh-eyes critique of an
earlier draft of this plan checked the diffusion/score-model literature specifically for
this: raw-scalar concatenation of a conditioning value is known to underperform a small
fixed sinusoidal (Fourier) embedding in that literature, though the effect is
depth-dependent --- FiLM/adaptive-norm modulation (the deeper diffusion-model-style fix) was
explicitly \emph{rejected} for v23 as disproportionate, since \code{VMFMixture}/\code{GMM1D}
are shallow (2-hidden-layer) heads where the signal-dilution problem FiLM solves barely
applies. The right-sized fix is the embedding alone:
\[
  e(\sigma) = \big[\sin(\sigma f_1), \cos(\sigma f_1), \dots, \sin(\sigma f_4), \cos(\sigma f_4)\big],
  \qquad f_j = \pi\!\cdot\!2^{\,j-1}\ (j=1,\dots,4),
\]
applied separately to $\log(\sigma_A+\sigma_0)$ and $\log(\sigma_n+\epsilon_n)$ (8 dims
each, 16 total), replacing the two raw scalars:
\[
  z' = \big[\,z \;;\; e(\log(\sigma_A+\sigma_0)) \;;\; e(\log(\sigma_n+\epsilon_n))\,\big],
  \qquad \sigma_0 = 1.0\text{\AA (matches \code{a100\_final}'s \code{SIGMA0\_A})}.
\]
$\sigma_n = \sigma_A/\text{scale}$ is blur normalized by the per-track scale
(\S\ref{sec:normalizer}); carry both the absolute and normalized forms, matching
\code{a100\_final}'s \code{PHYS\_LOG\_SIGMA\_A}/\code{PHYS\_LOG\_SIGMA\_N} convention, for the
same reasoning as before (absolute value reflects how much the normalizer itself was
corrected; normalized value reflects how much blur has washed out relative to track size).
\textbf{Zero-init the weight columns feeding these 16 embedded dims} in
\code{VMFMixture.net} and \code{GMM1D.net}'s first linear layer, so v23 at initialization is
behaviorally identical to v22 before any training signal distinguishes sigma values --- same
zero-init discipline as \code{a100\_final}, verified the same way (compare outputs with the
new columns zeroed vs.\ populated, at random init, confirm identical).

This is architecturally simpler to add here than it was for \code{a100\_final}'s
V-bottleneck design: there is no separate ``coefficient pathway'' vs.\ ``confidence
pathway'' to route sigma into selectively --- it is one MLP, and the embedded sigma is just
more input columns. $D_{\rm AUG}$ in \S\ref{sec:physfeat} is updated accordingly (16 dims
for sigma, not 2).

\subsubsection{Structural fallback: a zero-init residual on the raw-PCA axis}
\label{sec:fallback}
\textbf{Added after directly measuring \code{from\_scratch\_dr}'s real per-$\sigma$-tier
held-out performance} (a full 9-$\sigma\times$3-tier sweep run against its actual Stage-2
checkpoint, not the training log alone): despite a substantially \emph{better} zero-blur
starting point than R7 (4.9$^\circ$/7.8$^\circ$/11.7$^\circ$ at high/mid/low energy vs.\
R7's 8.8$^\circ$/17.8$^\circ$/27.8$^\circ$), \code{from\_scratch\_dr} collapses
\emph{faster and further} than R7 at every meaningful blur level in the target regime ---
e.g.\ high-energy $\sigma{=}50$nm: R7 32.8$^\circ$ vs.\ \code{from\_scratch\_dr}
\textbf{83.6$^\circ}$ (near-random); mid-energy $\sigma{=}10$nm: R7 39.7$^\circ$ vs.\
\textbf{82.0}$^\circ$. This holds even though \code{from\_scratch\_dr} was trained on
continuous blur from initialization, i.e.\ blur exposure alone (\S\ref{sec:training} item 1)
does not by itself prevent this collapse.

\textbf{Diagnosis.} R7's architecture cannot produce this failure mode \emph{by
construction}: its output is always some re-weighted blend of a handful of geometrically
robust candidates (PCA-type axes, attention centroids), so under distribution shift it can
at worst default toward the most robust candidate --- it has a structural floor. v22/v23's
unconstrained MLP head has no analogous floor: nothing prevents its output from behaving
arbitrarily once its input departs from what dominated training, independent of whether
$\sigma$-conditioning correctly informs it that this has happened. $\sigma$-conditioning
(\S\ref{sec:head} above) addresses whether the model \emph{knows} it should be less
confident; it does not by itself guarantee the point estimate \emph{degrades toward
something sensible} rather than arbitrarily. These are additive requirements, not
substitutes, exactly as \S\ref{sec:training} item 1 already argues for blur-exposure vs.\
sigma-conditioning.

\textbf{Mechanism, revives the highest-ranked unimplemented idea from the project's original
literature survey}: reparameterize \code{VMFMixture}'s mean direction as a zero-initialized
residual correction on top of the raw-PCA axis $u$ (already computed as part of
\S\ref{sec:physfeat}'s feature set, zero learned parameters), rather than an unconstrained
direct output:
\[
  \mu_k = \mathrm{normalize}\big(u + g(\sigma, z')\cdot \delta_k(z')\big),
\]
where $\delta_k(z')$ is the MLP's raw (previously final) output for component $k$, and
$g(\sigma,z')\in[0,1]$ is a learned, sigma-aware gate (a single sigmoid-activated scalar
per component, taking the embedded $\sigma$ features as input) that the network can use to
shrink the residual toward zero under high uncertainty. \textbf{Implementation note, precise
(external review flag --- there is no separate ``$\delta_k$ layer'' in the actual code to
zero as a unit):} \code{VMFMixture}'s \code{\_MLP} produces all $n_{\rm comp}{\times}5$
outputs (3 mu-dims $+$ $\kappa$ $+$ logit, interleaved per component) from one shared final
\code{nn.Linear} (\code{directional\_head.py} line~86). Zero only the weight/bias
\emph{rows} corresponding to the 3 mu-output dimensions per component, leaving
$\kappa$/logit rows at their normal initialization --- \emph{or}, verified equally safe,
zero the entire final layer (then $\kappa=\mathrm{softplus}(0)+\epsilon\approx0.70$ and
logits are uniform, both reasonable, non-degenerate starting values). Either is correct;
implement it as a row-level or full-layer zeroing on the existing shared linear layer, not as
a separate delta submodule, since none exists. This zero-init guarantees $\mu_k = u$ for
every component \emph{regardless of how the gate $g$ is initialized} --- if $\delta_k=0$
exactly, $g\cdot\delta_k=0$ for any value of $g$, so no special initialization of $g$ is
required for the guarantee to hold. v23 begins training
\emph{mathematically identical to plain PCA}, guaranteeing it can never do meaningfully worse
than PCA by construction, the same guarantee R7's candidate-vector-pool design provided
structurally. Unlike R7's design, this does not restrict the model's \emph{ceiling}: $g\to1$
and an expressive $\delta_k$ recovers the fully unconstrained v22 behavior wherever the data
supports it (i.e.\ at low blur, where v22's real strength lives), while $g\to0$ under high
blur recovers graceful PCA-like degradation instead of arbitrary collapse. This is a genuine
hedge, not a duplicate of $\sigma$-conditioning: it is possible for the gate mechanism to
still fail to learn the right shrinkage behavior from data alone, but the failure mode is now
bounded (worst case, an under-shrunk residual behaves like current v22/v23 without this
section) rather than unbounded (arbitrary collapse), which is a strictly safer starting
point.

\subsection{Physics-augmented scalar features --- v22's 5, plus higher-moment additions}
\label{sec:physfeat}
Keep v22's existing 5 features exactly (\code{src/utils/data\_utils.py} lines 1153--1167),
computed on the pre-truncation, pre-normalization cloud:
\[
  \phi = \big[\log(n_{\rm vac}{+}1),\ \log(\text{extent}{+}\epsilon),\
        \log(R_g{+}\epsilon),\ \lambda_1/\Sigma\lambda,\ \lambda_2/\Sigma\lambda \big]
\]
where $\lambda_1\ge\lambda_2\ge\lambda_3$ are the coordinate covariance eigenvalues.
\textbf{Note:} feature index 1 (``extent'') is the same L$_\infty$ statistic
\S\ref{sec:normalizer} replaces for the \emph{normalizer} --- keep it here as a
\emph{feature} unchanged (it is diagnostic information, not a scale to divide by, so its
known blur-drift is not a bug in this role), but be aware it is a redundant-but-differently-used
quantity; do not conflate the two roles when implementing.

\textbf{Add}: a scalar summary of the blur-corrected higher-moment (kurtosis-type) tensor
already derived and validated in \code{a100\_final} (\code{README.md} \S3, and
\code{validate\_math.py} TEST~B1/B2):
\[
  M = \tfrac{1}{n}\textstyle\sum_i p_i^2\, \mathrm{outer}(x_{c,i}, x_{c,i}),
  \qquad p_i = x_{c,i}\cdot u,\ u = \text{raw-PCA axis},
\]
with the exact blur correction
\[
  \hat A = M_{\rm obs} - s^2\big(\widehat{\langle q^2\rangle} I + 4\widehat{\langle q^2\rangle}\, u u^\top + \hat C\big)
           - s^4\big(I + 2uu^\top\big),
  \quad s^2 = \sigma_n^2\tfrac{n-1}{n},
\]
$\widehat{\langle q^2\rangle}=\langle p^2\rangle_{\rm obs}-s^2$, $\hat C = C_{\rm obs} - s^2 I$.
\code{a100\_final}'s own validation found the corrected tensor is decisively less biased
(bias reduction 2.6$\times$--196$\times$ across tested $\sigma$), \textbf{but} on individual
realizations the \emph{uncorrected} estimator's top eigenvector can outperform the corrected
one at high blur (uncorrected wins by 11$^\circ$ at $\sigma/R_g{=}2.37$ in
\code{a100\_final}'s test) because the uncorrected estimator has an accidental shrinkage
property toward the PCA axis that is itself beneficial there. \textbf{Do not pick one.} Feed
both as scalar features --- \textbf{using the signed cosine, not a raw angle}, per external
review: an eigenvector's sign is arbitrary (a sign flip near a degenerate case discontinuously
jumps a raw angle between values near $0^\circ$ and $180^\circ$, forcing the network to spend
capacity learning around a self-inflicted discontinuity). Sign-fix each eigenvector to $u$'s
hemisphere before use, or equivalently use the signed dot product directly:
\[
  c_{\rm corr} = \big|u\cdot \mathrm{eig}_1(\hat A)\big|, \qquad
  c_{\rm uncorr} = \big|u\cdot \mathrm{eig}_1(M_{\rm obs})\big|,
\]
(absolute value since axis direction, not sign, is what these features should convey ---
head/tail is handled separately, \S\ref{sec:loss}), plus the corresponding top eigenvalue
ratios as a rough confidence proxy. This lets the (now unconstrained, MLP-based) head learn
which one to trust as a function of $\sigma$ and energy, exactly the arbitration
\code{a100\_final} intended the reweighted candidate pool to provide, but without forcing a
static choice at feature-engineering time.

Total conditioning dimension: $D_{\rm AUG} = 384 + 5\,(\text{v22 phys}) + 4\,(\text{new
higher-moment scalars}) + 16\,(\sigma\text{-conditioning, sinusoidal-embedded}) = 409$.

\subsection{Normalizer --- replace with the blur-deconvolved radius of gyration}
\label{sec:normalizer}
Replace v22's legacy normalizer (\code{max\_extent = |centered|.max()}, an L$_\infty$
extreme-order-statistic proven to drift 2.27$\times$--7.77$\times$ under blur,
\code{data\_utils.py} line 1153) with \code{a100\_final}'s exact formula
(\code{data\_pipeline.py} lines 14--17, 64--70):
\[
  \text{scale} = \sqrt{\max\!\big(R^2_{\rm obs} - 3\sigma_A^2\tfrac{n-1}{n},\ \text{floor}\big)},
  \qquad \text{floor} = \max(\text{FLOOR\_ABS},\, (\text{FLOOR\_FRAC}\cdot\sigma_A)^2),
\]
where $R^2_{\rm obs} = \tfrac{1}{n}\sum_i |y_i-\bar y|^2$ is the observed mean-squared radius
of the blurred, centered cloud. This is unbiased in expectation
($\mathbb{E}[R^2_{\rm obs}] = R^2_{\rm true} + 3\sigma_A^2\tfrac{n-1}{n}$ exactly, no
approximation). Use the exact \code{FLOOR\_ABS}/\code{FLOOR\_FRAC} constants from
\code{a100\_final/data\_pipeline.py} (not re-derived here; read them directly from that file
at build time so v23 matches the already-validated constants rather than a plausible-sounding
guess).

\subsection{Graph construction --- replace fixed k=16 kNN with the blur-adaptive hybrid}
\label{sec:graph}
Replace v22's fixed \code{K\_NEIGHBORS=16} kNN with \code{a100\_final}'s hybrid
(\code{data\_pipeline.py} \S(5), \code{README.md} \S5): connect each point to all neighbors
within radius $r=\sqrt{H_0^2+\sigma_n^2}$ \emph{or} its \code{K\_MIN} nearest neighbors,
whichever gives more edges, capped at \code{K\_MAX}. Use \code{a100\_final}'s already-derived
constants: \textbf{$H_0=0.05$} (normalized units; chosen from the median unblurred
nearest-neighbor spacing across energy tiers --- low 0.100, mid 0.039, high 0.012 --- sitting
between the sparse low/mid-energy regime and the dense high-energy regime),
\textbf{K\_MIN=8} (anti-starvation floor for sparse low-energy tracks), \textbf{K\_MAX=24}
(hard cap, a documented approximation). Do not re-derive these; they were fit to the actual
data distribution already.

%==============================================================================
\section{Training procedure --- the frozen-backbone question, resolved}
\label{sec:training}
%==============================================================================

\textbf{Definitive recommendation: keep the two-stage structure (pretrain jointly, then
freeze the backbone and train a fresh head), but modify Stage 1 to include dynamic blur
exposure from the start.} Reasoning:

\begin{enumerate}[leftmargin=*]
\item v22's catastrophic collapse under blur has a confirmed, specific root cause distinct
      from sigma-blindness: \textbf{the backbone was never shown a single blurred example}
      (\code{smear\_sigma=0.0} default, never overridden in v22's \code{prepare\_data}). A
      frozen backbone that only ever learned to extract structure from perfectly clean point
      clouds cannot be expected to extract meaningful structure from a genuinely blurry one,
      \emph{regardless of whether the downstream head is told the blur level}. Sigma-conditioning
      tells the head how much to trust its input; it cannot manufacture structure the backbone
      never learned to see. These are two separate, additive requirements, not substitutes for
      each other.
\item \code{dr\_sigma} and \code{from\_scratch\_dr} already tested ``give the backbone real
      blur exposure'' (via unfreezing, and via training from scratch under blur respectively)
      and both failed --- \textbf{but neither included sigma-conditioning.} This session's
      established discipline is to change one variable at a time; unfreezing and
      training-from-scratch are both larger, more disruptive changes to the training
      procedure than adding blur exposure to an otherwise-unchanged Stage~1. The more
      surgical, single-new-variable test is: \emph{keep Stage 1's joint-training and
      freeze-then-fresh-head structure exactly as v22 built it, only add dynamic blur
      sampling to the data pipeline Stage 1 already uses} (the same \code{P\_ZERO},
      \code{MIN\_SIGMA\_A}, \code{MAX\_SIGMA\_A} continuous-sigma convention already used in
      R7/\code{a100\_final}), \emph{and} add sigma-conditioning to the head
      (\S\ref{sec:head}). This has not been tested in this combination.
\item If Stage 1 with blur exposure $+$ sigma-conditioned head still fails to generalize
      across sigma, \emph{then} escalate to unfreezing (re-testing \code{dr\_sigma}'s
      hypothesis, but this time with sigma-conditioning present) as a follow-up experiment,
      not a simultaneous one.
\end{enumerate}

\textbf{Stage 2} continues to freeze the (now blur-exposed) backbone and train a fresh
\code{DirectionalPosterior} from scratch, as v22 already does --- but Stage 2's training data
must \emph{also} be drawn under the same dynamic-blur distribution as Stage 1 (v22 currently
trains Stage 2 on the same clean Pool-A-only data as everything else); a head trained only on
clean examples, even sitting on top of a blur-aware backbone, would not itself learn to
modulate its output correctly across sigma.

\textbf{Data corpus:} train on \code{a100\_final}'s expanded \textbf{1{,}009{,}384-track v2
corpus} (24.9\%/30.0\%/45.2\% low/mid/high energy, deliberately skewed toward the
fixable mid/high-energy regime), not v22's original $\sim$320k-track pool
(\code{MAX\_TRAIN\_TRACKS=320\_000}, line 152) --- more data specifically benefits the
harder, blur-conditioned learning problem v23 takes on, and the pool is already built and
validated.

%==============================================================================
\section{Loss function}
\label{sec:loss}
%==============================================================================

\subsection{v22's actual mechanism, precisely}
The deployed posterior's own loss, \code{flow.loss(theta, z)}
(\code{DirectionalPosterior.loss}, \code{directional\_head.py} lines 158--163), is a fully
\textbf{signed} vMF-mixture NLL from the first training step --- there is no sign-blind
option inside \code{VMFMixture} itself. The curriculum described in prior session notes
(``sign-blind for \code{AUX\_SIGN\_WARMUP=30} epochs, then signed'') applies instead to a
\emph{separate, auxiliary} \code{DirectionHead} (\code{mu\_hat, kappa}) whose loss
(\code{axis\_aware\_vmf\_nll} for epoch $<$30, then \code{vmf\_nll}, \code{train\_gvp\_egnn\_v22.py}
lines 617--620) is added to the total Stage-1 objective purely to inject a stabilizing
gradient into the \emph{shared backbone} early in training, before the backbone's features are
reliable enough to support learning head/tail directly. Stage 2's own auxiliary loss
(lines 1222, 1261) stays \code{axis\_aware} throughout in the current v22 code — confirm
this is intentional (a design choice to keep Stage 2's auxiliary signal conservative since
Stage 1 already transferred head/tail-capable features) rather than an oversight before
carrying it forward unchanged.

\subsection{Interaction with $\sigma$-conditioning}
\textbf{Recommendation: do not make the curriculum's timing depend on $\sigma$.} The
curriculum exists to protect early-training gradient stability while backbone features are
generically unreliable; this is a training-dynamics concern, not a physics one, and there is
no evidence connecting it to blur specifically. Keep \code{AUX\_SIGN\_WARMUP=30} as-is. What
\emph{should} change: since Stage 1 now trains under a distribution of $\sigma$
(\S\ref{sec:training}), the warmup should be measured in a way that guarantees the
low-$\sigma$, most head/tail-informative examples have been seen a sufficient number of times
by epoch 30, not just that 30 epochs of wall-clock training have elapsed --- verify this
empirically rather than assume it transfers unchanged from the clean-data-only regime v22 was
tuned in.

%==============================================================================
\section{Calibration-specific design}
\label{sec:calibration}
%==============================================================================

The \code{a100\_final} finding motivating this section: ECE was \emph{worst} exactly where
the model was most accurate/confident (high energy, low blur, ECE $\sim$40\%, vs.\
$\sim$28--30\% at high blur). A proper scoring rule (vMF NLL) is a necessary but evidently
\emph{not sufficient} condition for calibration: with a flexible enough head and finite data,
gradient descent on NLL has every incentive to push $\kappa$ arbitrarily high wherever
point estimates are frequently close to correct, because NLL improves as $\kappa\to\infty$
for a correctly-placed mode --- this overfits to per-track noise in exactly the easy regime,
producing systematic overconfidence that does not generalize to held-out data.

\textbf{Definitive recommendation, two complementary mechanisms:}

\begin{enumerate}[leftmargin=*]
\item \textbf{An explicit concentration regularizer added directly to the training loss},
      analogous to the ``confidence penalty'' used to calibrate over-confident softmax
      classifiers (Pereyra et al.\ 2017): add a small penalty term proportional to the
      \emph{negative differential entropy} of the vMF mixture (a smooth, monotonic function
      of $\kappa$ for each component), \[ \mathcal{L}_{\rm calib} = \lambda \sum_k w_k\,
      \kappa_k, \] a simple, cheap, directly-tunable proxy for the same effect (entropy of a
      3-D vMF is monotonically decreasing in $\kappa$, so penalizing $\kappa$ directly is a
      valid, simpler stand-in for a full entropy term and avoids needing the vMF entropy's
      closed form). This applies uniformly across the training set — including easy,
      low-$\sigma$ examples — directly counteracting the runaway-confidence-when-accurate
      mechanism, without touching held-out data at all. $\lambda$ is a new hyperparameter;
      start small (e.g.\ tuned so $\mathcal{L}_{\rm calib}$ is $\sim$1\% of the direction NLL
      at initialization) and treat its selection as part of the Stage~2 build-and-validate
      gate (\S\ref{sec:eval}), not a guess frozen in advance.
\item \textbf{Calibration-aware checkpoint selection, not calibration-aware training.} Stage
      2 currently selects its best checkpoint by validation NLL (\code{best\_val\_loss\_s2}).
      Change the selection criterion to include validation-split coverage/ECE alongside NLL
      (e.g.\ select the checkpoint minimizing validation ECE among those within some small
      tolerance of the best validation NLL). \textbf{This is not a recalibration map} --- no
      additional trainable parameters are fit on validation or eval data, and no
      transformation is applied to the model's outputs after the fact. It only changes
      \emph{which already-trained epoch's weights are kept}, using a metric measured on data
      the model did not train on, which is a legitimate, standard form of model selection
      and is explicitly distinct from the post-hoc temperature scaling this document rules
      out.
\end{enumerate}

Report both the pre- and post-$\lambda$ (item 1 present vs.\ absent) coverage curves in the
build-and-validate gate so the regularizer's actual effect is measured, not assumed.

\textbf{Before trusting this mechanism's premise, verify it against real data (external
review flag, cheap, do first):} the runaway-$\kappa$-when-accurate explanation above is
argued from first principles, not yet checked against \code{a100\_final}'s actual $\kappa$
outputs. Pull the \code{a100\_final} ``small'' checkpoint and inspect mean $\kappa$ by
energy/$\sigma$ cell before finalizing $\lambda$'s design --- this five-minute check would
confirm or correct whether raw $\kappa$ growth is really the driver, versus e.g.\
mixture-weight collapse onto a single overconfident component, which would call for a
different regularizer term entirely.

\textbf{Named escalation path if $\lambda$ is insufficient:} the grid/Implicit-PDF-style
non-parametric posterior (built earlier this session, never deployed, kept only as an
untested fallback) is the correct next step specifically for calibration --- it directly
addresses a too-rigid posterior \emph{family}, which a $\kappa$ penalty on a fixed vMF
mixture cannot fix if the real problem is representational, not just a magnitude issue.

\textbf{Complementary, near-free addition, unconditional on whether $\lambda$ works:} train
2--3 seeds and ensemble (average) their posteriors. This is a well-established, essentially
free calibration improvement that does not require any architecture change and stacks with
either the $\kappa$-regularizer or the grid-posterior escalation --- include it in the build
regardless of how items 1--2 above perform.

%==============================================================================
\section{Evaluation methodology}
\label{sec:eval}
%==============================================================================

Use the identical convention already validated in \code{a100\_final/eval.py} and the
reference \code{full\_sweep\_9tier\_6nm\_EMA.csv}: the full \textbf{9-$\sigma$
$\times$ 3-energy-tier grid} ($\sigma\in\{0,1,3,6,10,20,30,50,100\}$nm $\times$
\{low, mid, high\}), each cell scored against the classical PCA/geometric baseline computed
with the same convention, 1000 held-out tracks/cell minimum. Calibration measured via the
same angular-coverage methodology, \textbf{reported per-$\sigma$-tier, not just pooled} ---
the whole reason the overconfidence-when-accurate pattern was caught in \code{a100\_final}
is that it was checked at multiple sigma levels rather than trusting one pooled number.

\textbf{Pre-registered success cells:} median axis error at mid-energy
$\sigma\in\{6,10,20\}$ and high-energy $\sigma\in\{6,10,20,30,50\}$ must improve over the
current best reference with a confidence interval excluding zero. \textbf{Status of that
reference, stated honestly:} \code{from\_scratch\_dr}'s real per-cell curve has now been
measured (full 9-$\sigma\times$3-tier sweep against its actual checkpoint, cited in
\S\ref{sec:fallback}) and is real, usable data. \textbf{v22's own zero-shot-under-blur curve
has not} --- only a single data point survives (77.8$^\circ$/57\% at 1nm, from
\code{dr\_sigma}'s docstring). \textbf{Run v22's actual \code{ROT\_FINAL} checkpoint through
the same full 9-$\sigma\times$3-tier sweep before training v23} (a cheap eval-only pass, no
training) so both real baselines --- v22 zero-shot and \code{from\_scratch\_dr} --- exist
before the success gate is evaluated against them, rather than comparing v23 against a
mixture of one measured curve and one single point.

\textbf{Pre-registered guard cells (hard stop if violated):} axis error and head-tail
accuracy at $\sigma\in\{0,1\}$ across all three tiers must not regress below v22's own
zero-blur numbers (6.8$^\circ$/97.0\% pooled) by more than a small, explicitly-agreed margin;
low-energy axis error at any $\sigma$ must not regress relative to the already-established
information-floor-saturated baseline (regression here is meaningless improvement-wise but a
regression would indicate something else broke).

\textbf{Calibration success:} mean ECE across the 27-cell grid materially below both v22's
zero-shot-under-blur collapse and \code{a100\_final}'s $\sim$30--40\%, \emph{and} no cell
should reproduce the ``ECE worst where most accurate'' inversion pattern found in
\code{a100\_final} --- check this explicitly as a named diagnostic, not just the aggregate
number.

%==============================================================================
\section{Staged build plan}
%==============================================================================

\begin{enumerate}[leftmargin=*]
\item \textbf{Component swap-in, zero-init verified (no training):} drop the V-bottleneck,
      confirm \code{DirectionalPosterior} runs unmodified; add the $\sigma$-conditioning
      inputs and the 4 new higher-moment scalar features with zero-init weights; swap in the
      \code{a100\_final} normalizer and blur-adaptive graph; add the $\kappa$-regularizer
      term (inactive, $\lambda=0$, to confirm it does not change forward-pass behavior before
      being turned on). Verify bit-identical outputs to plain v22 wherever the new inputs are
      zeroed, the same zero-init discipline used throughout \code{a100\_final}.
\item \textbf{Stage 1 modification, smoke-tested:} add dynamic blur sampling to Stage 1's
      data pipeline; confirm a short run trains without NaNs/crashes and that $\sigma{=}0$
      performance is not catastrophically different from v22's own smoke test.
\item \textbf{Full Stage 1 + Stage 2 training run} on the v2 corpus (1{,}009{,}384 tracks),
      with $\lambda>0$ active, on whatever GPU budget is available next (this is a smaller
      model than the \code{a100\_final} ``large'' config --- hidden=96/6 layers --- so it
      should be materially cheaper than the 30-hour A100 budget already spent; scope the
      actual budget once Stage 1's real throughput is measured, not guessed).
\item \textbf{Full evaluation} against the pre-registered grid, gates, and calibration
      diagnostics above.
\item \textbf{Only if the sigma-conditioned, blur-exposed two-stage model still fails the
      success gates:} escalate to testing unfreezing (re-run \code{dr\_sigma}'s hypothesis,
      now with sigma-conditioning present) as a distinct, separately-gated follow-up, not a
      simultaneous change.
\end{enumerate}

%==============================================================================
\section{Explicitly out of scope}
%==============================================================================
Restated for a reader who starts here: \textbf{v23 is not attempting to beat PCA at low
energy under real blur.} That regime's ceiling was established this session by three
independent methods and is treated as a physical fact, not an engineering target. v23's
scope is (a) recovering v22's proven zero-blur strength without regression, and (b)
achieving genuine blur-robustness at mid/high energy where the physics allows it, with
calibration that holds without any post-hoc correction.

\end{document}
